% =====================================================================
%  Chapter 26 opener -- "Chapter at a Glance", three reading paths + the
%  ladder staircase, THE running example (the shared 10-node rooted tree T
%  plus the recurring marked set S={4,8,6,10} and its virtual tree), the
%  who-owns-what comparison against Ch.16 (LCA), Ch.5 (Euler/tin), Ch.13
%  (tree DP), Ch.25 (link-cut) and Ch.22 (HLD), the recognition guide, a
%  constraint->technique table, and notation.
%  Exports fig:c30-tree, fig:c30-taxonomy.
%  No \term{lca} / \term{tin} here -- they are indexed by Ch.16/Ch.5; \term the
%  NEW objects "virtual tree", "auxiliary tree", "marked set", "compressed
%  edge" only.
% =====================================================================

\begin{chaptertheorybox}{Chapter at a Glance}
Chapter~16 (\chref{c11}) gave you \(\lca\): one DFS fixes \(\tin\), \(\tout\),
\(\depth\), and a binary-lifting table turns every ``lowest common ancestor'' and
``distance between two nodes'' into \(O(\log n)\) arithmetic
(\(\thmref{c11-distf}\)).  It closed by pointing forward -- ``a \emph{virtual tree}
(Ch.~26) compresses a \(k\)-node set via sort-by-tin plus the LCA of adjacent
pairs; this chapter's distance formula and lifting are the inner loop.''  This
chapter cashes that cheque.

A virtual tree answers a question no full tree DP can afford: you have one
\emph{fixed} tree on \(n\) nodes and a \emph{small set of \(k\) marked nodes}
(a colour class, a set of query terminals, a set of special cities), and you must
run some DP \emph{over just those \(k\) nodes and the paths between them} -- once
per colour, per query, per subset.  Doing it on the real tree is \(O(n)\) each
time; the virtual tree collapses the whole structure to at most \(2k-1\) nodes in
\(O(k\log k)\) -- the marked nodes plus the LCAs that connect them -- while an edge
keeps the \emph{length} of the path it swallowed, so every distance, subtree-size
and connectivity fact among the marked nodes survives untouched.

\textbf{You will learn to}
\begin{itemize}[nosep]
  \item build the \term{virtual tree} of a \term{marked set}: sort the \(k\) nodes by
        \(\tin\), insert \(\lca\) of each adjacent pair, dedup, then link every node
        to its nearest kept ancestor with a monotone stack -- \(\bigO{k\log k}\) and
        at most \(2k-1\) nodes (\(\secref{c30-construction}\), \thmref{c30-size},
        solved CF 2127E \emph{Ancient Tree}); \hfill\lv{3}--\lv{4}
  \item run a DP on the compressed tree, charging each \term{compressed edge} its real
        length \(\dist(u,v)=\depth(v)-\depth(u)\) and carrying the count of swallowed
        nodes, so the answer equals the \(O(n)\) DP done honestly (\(\secref{c30-dp}\),
        solved CF 2252F \emph{Spectral Components}); \hfill\lv{5}--\lv{6}
  \item spend the machine on the two workhorse families -- the sum of pairwise
        distances over a marked set (each edge contributes \(w\cdot\mathit{cnt}\cdot
        (k-\mathit{cnt})\)) and Steiner-like connected-subgraph counts
        (\(\secref{c30-apps}\)); \hfill\lv{5}
  \item read a statement and recognise the ``big tree, tiny marked set, re-run per
        colour/subset'' signature (\(\secref{c30-recognition}\)); \hfill\lv{4}
  \item know when a virtual tree is the WRONG tool -- bitmask Steiner DP, whole-tree
        DP, or a topology that changes (\(\secref{c30-boundary}\), the pick). \hfill\lv{5}
\end{itemize}

\textbf{Prerequisites.} Chapter~16 (\chref{c11}) -- \(\lca\), \(\dist\), binary
lifting: the inner loop, used and never re-derived.  Chapter~5 (\chref{c24}) -- the
\(\tin/\tout\) Euler order that is the sort key.  Chapter~13 (\chref{c10}) -- the
bottom-up subtree fold you re-run on the compressed tree.  Chapter~12 (\chref{c09})
-- \(\depth\) as the currency a compressed edge spends.  Chapter~25 (\chref{c23})
and Chapter~22 (\chref{c16}) -- the neighbours on either side: a link-cut tree is
this idea with a \emph{moving} topology; HLD is a fixed decomposition that does not
shrink with \(k\).
\end{chaptertheorybox}

% ---------------- Reading paths + ladder position ----------------
\begin{chaptertheorybox}{How to Read This Chapter (three paths)}
Level badges \(\lv3\)--\(\lv6\): \(\lv3\) is hand-building the virtual tree of a
small marked set and reading off its \(2k-1\) size; \(\lv{4}\) is the monotone-stack
construction coded and the \(\bigO{k\log k}\) bound justified; \(\lv{5}\) is a DP
on the compressed tree that charges compressed-edge lengths and swallowed-node
counts, or the pairwise-distance fold; \(\lv{6}\) is the Summit -- a virtual tree
over factorial nodes in an infinite divisor tree, solved by a weighted centroid.

\textbf{Construction path (\(\lv3\)--\(\lv{5}\), ``make the tree shrink'').}
\secref{c30-construction} in full -- sort-by-tin, adjacent LCAs, the stack,
\thmref{c30-size}, solved CF 2127E -- then \secref{c30-dp}, where the compressed
tree is fed a DP (solved CF 2252F), capped by the Summit \(\probref{c30-summit}\).

\textbf{Full-power path (\(\lv{4}\)--\(\lv{6}\), ``a DP over a \(k\)-set on a
\(10^5\) tree'').} \secref{c30-dp} for edge weights and the skipped-count carry,
\secref{c30-apps} for the distance-sum and connected-count families, then the
Summit, where the marked set is the factorial nodes and the payoff is a weighted
1-median.

\textbf{Boundary path (\(\lv{4}\)--\(\lv{5}\), ``is this even a virtual-tree job?'').}
\secref{c30-recognition} and \secref{c30-boundary}: the tell (tiny \(k\), re-run per
colour) and the four look-alikes that are not -- bitmask Steiner DP, whole-tree DP,
a moving forest for the link-cut tree (\chref{c23}), a static range for HLD
(\chref{c16}).
\end{chaptertheorybox}

\begin{figure}[H]
\centering
\begin{tikzpicture}[diagram, x=8.5mm, y=6.5mm]
  \foreach \st/\lab in {1/lv1,2/lv2,3/lv3,4/lv4,5/lv5,6/lv6} {
    \fill[\lab!22, draw=\lab, line width=0.8pt] (\st-1,0) rectangle (\st,\st*0.52);
    \node[font=\scriptsize\bfseries, text=\lab] at (\st-0.5,\st*0.52-0.42) {L\st};}
  \node[annot, anchor=south west, font=\scriptsize]
    at (0.05,0.62) {prereq (Ch.16): LCA, dist, tin -- the inner loop};
  \node[annot, anchor=south west, font=\scriptsize]
    at (1.05,1.14) {26.1 sort by tin + adjacent LCAs; dedup; the 2k-1 bound};
  \node[annot, anchor=south west, font=\scriptsize]
    at (2.05,1.66) {26.1 stack build, solved CF 2127E (per-colour VT)};
  \node[annot, anchor=south west, font=\scriptsize]
    at (3.05,2.18) {26.2 compressed-edge weight, skipped counts, solved CF 2252F};
  \node[annot, anchor=south west, font=\scriptsize, text=lv5]
    at (4.05,2.70) {26.3 distance-sum / connected-count; 26.4 recognition};
  \node[annot, anchor=south west, font=\scriptsize, text=lv6]
    at (5.05,3.22) {26.5 the pick; Summit CF 1292D (factorial VT + weighted centroid)};
\end{tikzpicture}
\caption{Chapter~26 on the six-step ladder: hand-building a virtual tree and
reading its size is \(\lv3\); the coded monotone-stack build and the \(O(k\log k)\)
bound are \(\lv{4}\); a DP that charges compressed-edge lengths and swallowed-node
counts (or the pairwise-distance fold) is \(\lv{5}\); the Summit (a virtual tree over
factorial nodes solved by a weighted centroid) is step~6.}
\end{figure}

% ---------------- The running example ----------------
\subsection*{The Running Example: One Fixed Tree, One Tiny Marked Set}

Every section in this chapter runs on the same rooted tree \(T\) on \(n=10\)
nodes, fixed once here and authoritative (if a later trace disagrees, this figure
wins).  It is a single tree rooted at \(1\); parent-to-child edges, 1-based:
\[
  1\!\to\!2,\quad 1\!\to\!3,\quad 2\!\to\!4,\quad 2\!\to\!5,\quad
  3\!\to\!6,\quad 3\!\to\!7,\quad 5\!\to\!8,\quad 5\!\to\!9,\quad 7\!\to\!10 .
\]
A left-to-right DFS assigns entry times \(\tin\): node \(1\!:\!1\), \(2\!:\!2\),
\(4\!:\!3\), \(5\!:\!4\), \(8\!:\!5\), \(9\!:\!6\), \(3\!:\!7\), \(6\!:\!8\),
\(7\!:\!9\), \(10\!:\!10\).  Depths are \(\depth(1)=0\);
\(\depth(2)=\depth(3)=1\); \(\depth(4)=\depth(5)=\depth(6)=\depth(7)=2\);
\(\depth(8)=\depth(9)=\depth(10)=3\).  The recurring \term{marked set} is
\(S=\{4,8,6,10\}\) (\(k=4\)) -- think of them as the four special vertices a query
cares about.  Their virtual tree \(A(S)\) keeps \(S\) plus the LCAs of adjacent
pairs in \(\tin\)-order, \(\{\lca(4,8),\lca(8,6),\lca(6,10)\}=\{2,1,3\}\), giving
\(7=2k-1\) nodes and six \term{compressed edges}; nodes \(5,7,9\) vanish into edge
lengths (\(\secref{c30-construction}\)).

\begin{figure}[H]
\centering
\begin{tikzpicture}[diagram, scale=0.9, transform shape,
   tnode/.style={circle, draw=primaryblue!70, fill=primaryblue!8, line width=0.7pt,
                 minimum size=6.2mm, font=\scriptsize\bfseries, inner sep=1pt},
   mk/.style={circle, draw=warningred, fill=warningred!16, line width=1.1pt,
             minimum size=6.8mm, font=\footnotesize\bfseries, inner sep=1pt},
   kee/.style={circle, draw=successgreen!75, fill=successgreen!12, line width=0.9pt,
             minimum size=6.2mm, font=\scriptsize\bfseries, inner sep=1pt},
   drop/.style={circle, draw=gray!60, fill=gray!8, line width=0.5pt, dashed,
             minimum size=5.4mm, font=\scriptsize\footnotesize, inner sep=1pt, text=gray}]
  % left: the real tree T
  \node[mk] (t1) at (0,0) {1};
  \node[kee] (t2) at (-2.2,-1.3) {2};
  \node[kee] (t3) at (2.2,-1.3) {3};
  \node[mk] (t4) at (-3.4,-2.7) {4};
  \node[drop] (t5) at (-1.2,-2.7) {5};
  \node[mk] (t6) at (1.5,-2.7) {6};
  \node[drop] (t7) at (3.3,-2.7) {7};
  \node[mk] (t8) at (-2.0,-4.1) {8};
  \node[drop] (t9) at (-0.5,-4.1) {9};
  \node[mk] (t10) at (3.3,-4.1) {10};
  \draw (t1)--(t2) (t1)--(t3) (t2)--(t4) (t2)--(t5) (t3)--(t6) (t3)--(t7)
             (t5)--(t8) (t5)--(t9) (t7)--(t10);
  \node[annot, font=\scriptsize, anchor=south] at (0.4,0.35) {the fixed tree T (n=10)};
  % right: the virtual tree
  \node[kee] (v1) at (8.2,0) {1};
  \node[kee] (v2) at (6.6,-1.6) {2};
  \node[kee] (v3) at (9.8,-1.6) {3};
  \node[mk] (v4) at (5.6,-3.3) {4};
  \node[mk] (v8) at (7.4,-3.3) {8};
  \node[mk] (v6) at (9.0,-3.3) {6};
  \node[mk] (v10) at (10.8,-3.3) {10};
  \draw[very thick] (v1)--node[left,font=\scriptsize,text=gray]{1} (v2)
        (v1)--node[right,font=\scriptsize,text=gray]{1} (v3)
        (v2)--node[left,font=\scriptsize,text=gray]{1} (v4)
        (v2)--node[right,font=\scriptsize,text=gray]{2} (v8)
        (v3)--node[left,font=\scriptsize,text=gray]{1} (v6)
        (v3)--node[right,font=\scriptsize,text=gray]{2} (v10);
  \node[annot, font=\scriptsize, anchor=south] at (8.4,0.35) {virtual tree A(S), 2k-1=7 nodes};
\end{tikzpicture}
\caption{The running example.  Left: the fixed tree \(T\); red nodes are the marked
set \(S=\{4,8,6,10\}\), green nodes are the LCAs the construction must add
(\(\{1,2,3\}\)), grey dashed nodes \(\{5,7,9\}\) are \emph{swallowed}.  Right: the
virtual tree \(A(S)\) -- every real edge between two kept nodes is replaced by a
\term{compressed edge} (thick) carrying the number of real edges it swallowed
(\(2\!\to\!8\) has length \(2\) along \(2\text{-}5\text{-}8\)).  The distance and
ancestor relation between any two red/green nodes is exactly the real tree's.}
\label{fig:c30-tree}
\end{figure}

% ---------------- Which job, cost chart ----------------
\subsection*{Which Job? (the question-shape chart)}

\begin{center}
\small
\renewcommand{\arraystretch}{1.25}
\setlength{\tabcolsep}{3pt}
\begin{tabular}{@{}p{0.31\linewidth} p{0.17\linewidth} p{0.28\linewidth} >{\raggedright\arraybackslash}p{0.14\linewidth}@{}}
\toprule
\textbf{Question} & \textbf{Technique} & \textbf{Why not a full-tree DP} & \textbf{Where} \\
\midrule
a DP over \(k\) marked nodes, re-run per colour/query & virtual tree + DP & \(O(n)\) per copy blows the budget; VT is \(O(k\log k)\) & \secref{c30-construction} \\
sum of \(\dist\) over all pairs of a \(k\)-set & VT, edge \(\cdot\)cnt\(\cdot(k\)-cnt) & pair walk is \(O(kn)\); VT touches \(2k\) edges & \secref{c30-apps} \\
cheapest connected subgraph spanning \(k\) terminals & VT + fold on it & the terminals' Steiner tree has \(O(k)\) branching nodes & \secref{c30-dp} \\
for a set of nodes, is \(u\) an ancestor of \(v\)? & \(\tin/\tout\) test & free once \(\tin\) is known (Ch.5) & \secref{c30-construction} \\
terminals joined through \emph{new} branch nodes & \emph{bitmask Steiner DP} & VT only keeps existing nodes, cannot invent terminals & \secref{c30-boundary} \\
the tree \emph{changes} between queries & \emph{go to link-cut} & VT assumes a fixed tree; edits need Ch.25 & \chref{c23} \\
\bottomrule
\end{tabular}
\end{center}

\begin{insightbox}[One promise decides everything: keep the nodes that branch]
A virtual tree keeps exactly the marked nodes and the nodes where the paths
between them \emph{branch} -- and a branch point of marked nodes is always an LCA of
two of them.  Sorting the \(k\) marked nodes by \(\tin\) makes every needed LCA the
LCA of two \emph{adjacent} marked nodes, so \(k-1\) LCA calls (plus dedup) already
close the set under ancestry.  Compress every maximal chain of unmarked, non-branch
nodes into a single weighted edge, and you get a tree with at most \(2k-1\) nodes
that answers every distance, ancestor and connectivity question among the marked
set \emph{exactly} -- \thmref{c30-size}.  The whole chapter is one idea: \emph{make
the tree shrink to \(O(k)\) without changing what the \(k\) nodes know about each
other.}
\end{insightbox}

% ---------------- Taxonomy chart ----------------
\begin{figure}[H]
\centering
\begin{tikzpicture}[diagram, scale=0.95, transform shape,
    mnode/.style={circle, draw=#1, line width=0.9pt, fill=#1!8, minimum size=9mm,
                  font=\scriptsize\bfseries, text=#1!60!black}]
  \node[mnode=primaryblue] (a) at (0,0)   {close};
  \node[annot, font=\scriptsize, anchor=west, text width=9.8cm, align=left] at (1.7,0) {sort the k marked nodes by tin; for each adjacent pair take lca and add it; dedup. This set A(S) is closed under lca, and has at most 2k-1 nodes (26.1)};
  \node[mnode=successgreen] (b) at (0,-1.1) {link};
  \node[annot, font=\scriptsize, anchor=west, text width=9.8cm, align=left] at (1.7,-1.1) {sweep A(S) by tin with a monotone stack: pop until the stack top is an ancestor of the current node, then make the current node a virtual child of the top. Edge weight = depth difference (26.1--26.2)};
  \node[mnode=difforange]   (c) at (0,-2.2) {fold};
  \node[annot, font=\scriptsize, anchor=west, text width=9.8cm, align=left] at (1.7,-2.2) {run the usual bottom-up tree DP on the compressed tree, but charge each virtual edge its real length and carry the count of swallowed nodes, so the value equals the honest O(n) DP (26.2, 26.3; solved CF 2252F)};
  \node[mnode=warningred]   (d) at (0,-3.3) {recognise};
  \node[annot, font=\scriptsize, anchor=west, text width=9.8cm, align=left] at (1.7,-3.3) {the tell is a fixed big tree, a small marked set, and a computation repeated per colour / per query / per subset with sum k << total budget (26.4)};
  \node[mnode=softpurple]   (e) at (0,-4.4) {reject};
  \node[annot, font=\scriptsize, anchor=west, text width=9.9cm, align=left] at (1.7,-4.4) {whole tree needed -> tree DP; terminals can meet at NEW nodes -> bitmask Steiner DP; topology moves -> link-cut (Ch.25); static ranges -> HLD/Euler (Ch.22/Ch.5) (26.5)};
\end{tikzpicture}
\caption{The chapter's five-state taxonomy.  \emph{close} turns the marked set into
an LCA-closed set; \emph{link} builds the compressed tree with a stack; \emph{fold}
runs a weighted DP on it; \emph{recognise} spots the signature in a statement;
\emph{reject} hands the job to the right neighbour.  Each solved problem sits on one
row: CF 2127E on \emph{close}+a count, CF 2252F on \emph{fold}, the Summit CF 1292D
on \emph{link}+a weighted centroid.}
\label{fig:c30-taxonomy}
\end{figure}

% ---------------- Who owns what ----------------
\subsection*{Who Owns What (so nothing is learned twice)}

\begin{center}
\small
\renewcommand{\arraystretch}{1.25}
\setlength{\tabcolsep}{3pt}
\begin{tabular}{@{}p{0.28\linewidth} p{0.24\linewidth} p{0.44\linewidth}@{}}
\toprule
\textbf{Idea} & \textbf{Owner} & \textbf{This chapter's share} \\
\midrule
\(\lca\), \(\dist\), binary lifting & Ch.~16 (\chref{c11}) & the inner loop of \emph{close} and \emph{link}; used, cited by \thmref{c11-distf}, never re-derived \\
\(\tin/\tout\) flatten, ancestor test & Ch.~5 (\chref{c24}) & the sort key and the stack's ``is ancestor'' predicate; never re-taught \\
bottom-up subtree fold & Ch.~13 (\chref{c10}) & the same fold, run on the \(O(k)\) compressed tree with edge weights (26.2) \\
\(\depth\) as a distance currency & Ch.~12 (\chref{c09}) & a compressed edge's weight is a \(\depth\) difference (26.1--26.2) \\
dynamic forest, moving topology & Ch.~25 (\chref{c23}) & explicitly out; 26.5 routes ``edges change'' to the link-cut tree \\
static path/subtree \emph{ranges} & Ch.~22 HLD / Ch.~5 Euler & explicitly out; those do not shrink with \(k\) and are the boundary's first rejects \\
marked-set queries via lifting & Ch.~16 Summit & 26.x reuses that set-LCA insight but builds the whole induced tree, not one range \\
\bottomrule
\end{tabular}
\end{center}

% ---------------- Recognition guide ----------------
\begin{recognitionbox}[Pattern card -- ``a big fixed tree, a tiny set I re-run on'']
\reckey{statements that fix one tree of \(n\le 2\cdot10^5\) nodes and then hand you
a \emph{small} set of special vertices -- one colour class, one query's terminal list,
one subset -- and ask you to compute something over exactly that set and the paths
linking it (a DP, a distance sum, a min connected subgraph, ``for each colour over its
vertices\ldots'').  The tell is that the same \(O(n)\) tree DP would be re-run
\emph{many times}, once per subset, and the total number of marked vertices across all
runs is small -- \(\sum k \ll n\cdot(\text{runs})\).}
\recuse{
  \item per-colour / per-subset DP over that subset's induced subtree -- virtual tree + fold (\(\secref{c30-dp}\));
  \item sum of distances over all pairs of a \(k\)-set -- virtual tree, per-edge \(w\cdot\mathit{cnt}\cdot(k-\mathit{cnt})\) (\(\secref{c30-apps}\));
  \item minimum connected subgraph spanning \(k\) marked nodes (Steiner \emph{on a tree}) -- virtual tree DP (\(\secref{c30-apps}\)).
}
\recnotuse{
  \item the answer needs every node's own value, not just the marked ones -- full tree DP / rerooting (\chref{c10}, \chref{c20});
  \item terminals may be joined at nodes that are \emph{not} in the tree you were given / new branch nodes -- bitmask Steiner DP, not a virtual tree (\(\secref{c30-boundary}\));
  \item the tree itself is edited between queries -- link-cut tree (\chref{c23});
  \item you query ranges on a \emph{fixed} tree and nothing shrinks with \(k\) -- HLD / Euler+Fenwick (\chref{c16}, \chref{c24}).
}
One-line test: \emph{one fixed tree, a marked set of size \(k\ll n\), a computation
repeated over that set -- then collapse to \(O(k)\) with a virtual tree; if every node
matters or the set never shrinks, keep the full-tree engine.}
\end{recognitionbox}

% ---------------- Constraint -> technique ----------------
\subsection*{Reading the Constraints (which technique the limits want)}

\begin{center}
\small
\renewcommand{\arraystretch}{1.25}
\begin{tabular}{@{}p{0.34\linewidth} p{0.40\linewidth} p{0.20\linewidth}@{}}
\toprule
\textbf{Constraint signature} & \textbf{What it rules in/out} & \textbf{Technique} \\
\midrule
fixed tree \(n\le2\!\cdot\!10^5\), and \(\sum(\text{marked})\) over all runs \(\le 2\!\cdot\!10^5\) & a full \(O(n)\) DP per run is too slow; each run wants \(O(k\log k)\) & \secref{c30-construction} \\
``for each colour / each query, over its \(k\) vertices\ldots'' & the object lives on the marked set's induced subtree, not the whole tree & \secref{c30-dp} \\
sum of distances over all pairs of a small set & \(O(k^2)\) pairs too slow; fold per virtual edge & \secref{c30-apps} \\
Steiner tree / min connected subgraph on a TREE (terminals are existing nodes) & the Steiner nodes are already the marked set + their LCAs & \secref{c30-apps} \\
terminals may meet at \emph{new} nodes, \(k\le18\) & a virtual tree cannot create terminals; need bitmask DP & \secref{c30-boundary} \\
\bottomrule
\end{tabular}
\end{center}

% ---------------- Notation ----------------
\subsection*{Notation Used Only Here}

The fixed tree is \(T=(V,E)\), \(|V|=n\), rooted at \(1\).  A \term{marked set}
\(S\subseteq V\), \(|S|=k\), is the small set of vertices a single run cares about.
The \term{virtual tree} \(A(S)\) is the minimal ancestor-closed set containing \(S\)
(equivalently \(S\) together with \(\lca(x,y)\) for adjacent \(x,y\) in \(\tin\)-order),
viewed as a tree whose edges are \term{compressed edges}: a virtual edge \((u,v)\),
\(u\) a proper ancestor of \(v\), stands for the whole real path \(u\!\to\!v\) and
carries weight \(\dist(u,v)=\depth(v)-\depth(u)\).  \(\mathit{cnt}(v)\) is the number
of marked nodes in \(v\)'s virtual subtree; \(\size(v)\) is the number of \emph{real}
tree nodes swallowed inside \(v\)'s virtual subtree (used when the DP must charge
skipped nodes).  \(\Omega(x)\) is the number of prime factors of \(x\) with
multiplicity -- it is the depth of node \(x\) in the Summit's divisor tree.
\(\depth,\tin,\tout,\lca,\dist\) are Ch.~5/Ch.~16's, used unchanged.  Every section
runs on the \(T\) and \(S=\{4,8,6,10\}\) above; if a table here and a trace there
disagree, this opener is authoritative.
